\section{Feature Engineering}\label{sec:feature_engineering}

\textit{Feature engineering} adalah proses membuat fitur-fitur (kolom data) baru atau mengubah fitur yang sudah ada agar representasi data menjadi lebih informatif bagi algoritma \textit{machine learning}. Menurut Kuhn dan Johnson (2019) dalam \textit{Feature Engineering and Selection: A Practical Approach for Predictive Models}, \textit{feature engineering} yang efektif sering kali menjadi faktor penentu utama keberhasilan sebuah model prediktif, dikarenakan algoritma tidak dapat menemukan pola yang kompleks jika data mentahnya tidak merepresentasikan masalah dengan baik. Tujuannya adalah membantu model untuk belajar dari data dan mengidentifikasi pola yang menyimpang dari kondisi normal (anomali).

Berdasarkan data operasional sistem penyewaan, \textit{feature engineering} berfokus pada ekstraksi fitur temporal (waktu) dan pemrosesan informasi kualitatif (teks) guna mendeteksi pola perilaku penyewa yang mencurigakan. Dalam konteks data penyewaan, berikut adalah contoh perhitungan \textit{feature engineering} dari fitur-fitur utama yang dikembangkan beserta landasan teorinya.

\begin{enumerate}
    \item \textbf{Indikator Pelanggaran Berdasarkan Komentar (\textit{Is Melanggar})} \\
    Dalam sistem pencatatan operasional, anomali sering kali tidak tercatat dalam bentuk metrik angka, melainkan melalui catatan bebas (\textit{free-text}) atau log laporan dari staf \cite{aggarwal2017outlier}. Jika kolom komentar pada riwayat penyewa tidak kosong (\textit{not null}), sistem akan menghasilkan fitur biner baru bernilai 1 (\textit{True}), dan 0 (\textit{False}) jika sebaliknya. Hal ini berfungsi sebagai sinyal kuat bagi algoritma mengenai riwayat pelanggaran tata tertib.

    \item \textbf{Durasi Sewa dalam Jam (\textit{Duration Hours})} \\
    Fitur ini menghitung total durasi penyewaan secara kontinu dalam satuan jam, berdasarkan selisih antara waktu \textit{check-out} dan waktu \textit{check-in}.
    
    {Contoh kasus perhitungan:}
    \begin{center}
        \begin{tabular}{|c|c|c|}
            \hline
            \textbf{Waktu Check-in} & \textbf{Waktu Check-out} \\
            \hline
            2025-05-18 20:00:00 & 2025-05-18 23:00:00\\
            \hline
            2025-05-19 14:00:00 & 2025-05-20 12:00:00\\
            \hline
        \end{tabular}
    \end{center}
    
    Persamaan umum yang digunakan adalah:
    \[
    \text{Duration Hours} = \frac{(\text{Check-out} - \text{Check-in})_{\text{detik}}}{3600}
    \]
    
    \begin{itemize}
        \item \textbf{Kasus 1 (Penyewaan di hari yang sama):} \\
        Selisih waktu (20:00:00 hingga 23:00:00) adalah 3 jam. Dalam satuan detik bernilai $3 \times 3600 = 10.800$ detik.
        \[
        \text{Durasi}_1 = \frac{10.800}{3600} = 3{,}0 \text{ jam}
        \]
        \item \textbf{Kasus 2 (Penyewaan melintasi hari):} \\
        Selisih waktu (19 Mei 14:00:00 hingga 20 Mei 12:00:00) adalah 22 jam. Dalam satuan detik bernilai $22 \times 3600 = 79.200$ detik.
        \[
        \text{Durasi}_2 = \frac{79.200}{3600} = 22{,}0 \text{ jam}
        \]
    \end{itemize}

    Fitur ini krusial untuk mendeteksi penyewa ``transit''. Baesens, Van Vlasselaer, dan Verbeke  \cite{baesens2015fraud} menyatakan bahwa ekstraksi fitur temporal adalah teknik paling efektif dalam sistem deteksi anomali. Sejalan dengan temuan tersebut, Gupta et al.  \cite{gupta2014outlier} menegaskan bahwa anomali temporal sering muncul dalam bentuk durasi kejadian yang terlampau singkat. Dalam konteks sewa properti, durasi inap yang sangat pendek (\textit{micro-stays}) berkorelasi positif dengan aktivitas ilegal.

    \item \textbf{Ekstraksi Waktu Check-in (\textit{Hour} \& \textit{Day of Week})} \\
    Fitur ini memecah informasi waktu kedatangan menjadi komponen jam operasional (0--23) dan indeks hari dalam seminggu (0 = Senin, 6 = Minggu). Perilaku manusia memiliki pola musiman dan temporal yang dapat diprediksi \cite{zheng2014urban}. Kedatangan pada jam tidak lazim (misalnya pukul 03.00 dini hari) dapat menjadi indikator awal (\textit{early warning}) dari profil anomali jika digabungkan dengan fitur durasi yang singkat.

    \item \textbf{Indikator Akhir Pekan (\textit{Is Weekend Check-in})} \\
    Fitur biner ini menandai apakah penyewaan dilakukan pada akhir pekan (Sabtu atau Minggu). Menurut Aggarwal  \cite{aggarwal2017outlier}, waktu merupakan atribut kontekstual (\textit{contextual attribute}) yang mendefinisikan batasan normalitas suatu data. Lonjakan penyewaan di akhir pekan adalah hal wajar, namun bisa menjadi sinyal mencurigakan jika terjadi di pertengahan minggu kerja. Chandola, Banerjee, dan Kumar  \cite{chandola2009anomaly} mendefinisikan fenomena ini sebagai anomali kontekstual. Fitur ini membantu model menghindari kesalahan klasifikasi (\textit{false positive}) dengan memberikan konteks waktu yang tepat.
\end{enumerate}

Selain melakukan ekstraksi fitur seperti waktu dan hari, tahap persiapan data untuk \textit{machine learning} harus mempertimbangkan rentang skala matematis dari data asli. Algoritma seperti \textit{Support Vector Machine} (SVM) sangat sensitif terhadap nilai fitur yang memiliki variasi skala berbeda-beda karena mengandalkan kalkulasi jarak numerik antar titik data. Oleh karena itu, fitur numerik harus distandardisasi. Menurut Bishop  \cite{bishop2006pattern} dalam \textit{Pattern Recognition and Machine Learning}, standardisasi (\textit{StandardScaler}) beroperasi dengan menyesuaikan setiap nilai fitur hingga memiliki rata-rata nol ($\mu = 0$) dan standar deviasi satu ($\sigma = 1$). Persamaan matematika standardisasi adalah $z = \frac{x - \mu}{\sigma}$, di mana fitur hasil standardisasi dipusatkan agar algoritma konvergen lebih cepat dan secara geometris proporsional tanpa ada satu fitur bernilai ribuan yang mendominasi fitur bernilai desimal kecil.

Selama mengeksekusi tahapan standardisasi maupun manipulasi fitur lainnya, kehati-hatian secara ilmiah harus dikedepankan guna mencegah fenomena kebocoran data (\textit{data leakage}). Kaufman et al.  \cite{kaufman2011leakage} mendefinisikan \textit{data leakage} sebagai suatu kelalaian di mana informasi krusial dari kumpulan data pengujian (\textit{testing set}) diam-diam masuk, membaur, atau ikut terhitung saat proses pembentukan model pada data pelatihan (\textit{training set}). Apabila misalnya nilai rata-rata keseluruhan (gabungan rasio data latih dan pengujian) dipakai pada proses \textit{StandardScaler}, model akan mengingat pola dari data masa depan, sehingga membuat akurasinya tampak fantastis secara keliru, namun akan gagal secara drastis saat membedakan kasus anomali baru karena model gagal belajar menggeneralisasi. 

% Berikut adalah contoh implementasi kode Python menggunakan pustaka \texttt{pandas} untuk mengeksekusi proses \textit{feature engineering} tersebut.

% \begin{lstlisting}[language=Python, caption={Implementasi Feature Engineering Menggunakan Python}]
% # Pastikan kolom yang dibutuhkan tersedia di dataframe
% cols_to_check = [
%     'tanggal_checkin', 'waktu_checkin', 'tanggal_checkout', 'waktu_checkout',
%     'lama_menginap', 'lantai', 'tower', 'status_kewarganegaraan', 
%     'metode_pembayaran', 'jenis_sewa', 'komentar' 
% ]

% for col in cols_to_check:
%     if col not in df.columns:
%         df[col] = np.nan
        
% # Fitur Indikator Pelanggaran (Berdasarkan Komentar)
% # Menghasilkan nilai 1 (True) jika komentar ada dan bukan string kosong
% df['melanggar'] = df['komentar'].notna() & (df['komentar'].str.strip() != '')
% df['melanggar'] = df['melanggar'].astype(int)

% # Fungsi Penggabungan Tanggal dan Waktu
% def combine_datetime(date_col, time_col):
%     dt_str = df[date_col].astype(str).str.strip() + ' ' + df[time_col].astype(str).str.strip()
%     return pd.to_datetime(dt_str, errors='coerce')

% df['checkin_dt'] = combine_datetime('tanggal_checkin', 'waktu_checkin')
% df['checkout_dt'] = combine_datetime('tanggal_checkout', 'waktu_checkout')

% # Perhitungan Durasi Sewa dalam Jam (Duration Hours)
% df['duration_hours'] = (df['checkout_dt'] - df['checkin_dt']).dt.total_seconds() / 3600
% # Mengisi nilai kosong dengan 0 dan mencegah nilai negatif (kesalahan input)
% df['duration_hours'] = df['duration_hours'].fillna(0).clip(lower=0)

% # Ekstraksi Fitur Waktu (Jam dan Hari)
% df['checkin_hour'] = df['checkin_dt'].dt.hour.fillna(0)       # Jam (0-23)
% df['checkout_hour'] = df['checkout_dt'].dt.hour.fillna(0)     # Jam (0-23)
% df['checkin_dow'] = df['checkin_dt'].dt.dayofweek.fillna(-1)  # Hari (0=Senin, 6=Minggu)

% # Fitur Biner Akhir Pekan
% df['is_weekend_checkin'] = df['checkin_dow'].isin([5, 6]).astype(int)

% # Normalisasi Tipe Data
% df['lama_menginap'] = pd.to_numeric(df['lama_menginap'], errors='coerce').fillna(0)
% \end{lstlisting}

% \textbf{Penjelasan Kode Program:}
% \begin{itemize}
%     \item \texttt{df['komentar'].notna()}: Fungsi untuk mendeteksi baris data yang memiliki entri kualitatif (bukan \textit{null} atau \textit{NaN}). Hasil logika \textit{boolean} kemudian dikonversi menjadi numerik melalui \texttt{astype(int)} agar kompatibel dengan model \textit{machine learning}.
%     \item \texttt{combine\_datetime()}: Menggabungkan variabel tanggal dan waktu yang awalnya terpisah (bertipe \textit{string}) menjadi satu objek tipe \texttt{datetime} terpadu untuk operasi matematika waktu.
%     \item \texttt{.dt.total\_seconds() / 3600}: Mengonversi nilai \textit{timedelta} (selisih waktu absolut) menjadi representasi numerik desimal dalam satuan jam (\textit{hours}).
%     \item \texttt{.clip(lower=0)}: Sebuah fungsi proteksi untuk membatasi batas bawah matriks menjadi nol, guna mencegah kemunculan durasi negatif akibat anomali kesalahan ketik saat staf memasukkan data \textit{check-in} dan \textit{check-out}.
% \end{itemize}
